\section{Robustness to Core Threats}\label{sec:robustness}

The robustness checks follow the empirical chain. The legal shock must line up
with the timing of prices; drop-out bids must recover the cost objects needed
for the order-statistic comparison; the post-policy SME pool must not be the
sole source of the result; coordination must not intensify after restriction;
and the preference benchmark must not be a mechanical artifact. The question is
whether any of these threats overturns the conclusion that bidder exclusion is
the main source of the price increase.

\subsection{Reduced-form timing and scale}

The reduced-form evidence is used for timing and scale, not for the structural
decomposition. The \structuralWindowM-month benchmark in
Table~\ref{tab:did_benchmark_v8} is \didPriceEighteenPbu{} in the paper's
DiD-in-reverse convention, implying a \didLowerPct--\didUpperPct{} percent
price movement. The estimate is stable across the 6-, 12-, and 18-month
windows. Excluding the BEC enablement months before the mass-take-up cutoff
changes the coefficient from \phasedDidFull{} to \phasedDidTrunc{} with a
standard error of \phasedDidSe, and placebo cutoffs before the reform are
smaller in magnitude (\placeboSeventeenSep, \placeboSeventeenMar, and
\placeboSeventeenJun) than the real-cutoff estimates.

The conservative reading remains necessary. Group~65 is not perfectly balanced
against the pooled controls: \balanceNVarsAboveTen{} of \balanceNVarsTotal{}
pre-period covariates exceed the 0.10 normalized-difference threshold, and
\balanceNVarsAboveTwentyFive{} exceed 0.25. Modern DiD variants also reduce
precision: the BJS imputation estimate is \didAltBjsEst{} and the
Callaway--Sant'Anna estimate is \didAltCsEst. The reduced form therefore
validates the episode's sign and timing; it does not identify the exclusion and
protected-pool channels. Online Appendix~OA-B reports the balance table,
placebos, and alternative estimators.

\subsection{Cost recovery and auction primitives}

The cost-recovery concern is that the decomposition could be an artifact of
winner censoring, common auction-level shocks, or bidder dependence. The
diagnostics do not support that reading.
Alternative treatments of winner censoring give similar net price effects:
\costRegimeLosersNp/\costRegimeAllNp/\costRegimeTurnNp{} in non-pharmaceuticals
under losers-only/all-bidders/Turnbull inputs, and
\costRegimeLosersPh/\costRegimeAllPh/\costRegimeTurnPh{} in pharmaceuticals.
The Turnbull within-auction share remains large
(\turnbullShareNp{} percent in non-pharmaceuticals and \turnbullSharePh{}
percent in pharmaceuticals), so the result is not driven by treating the
winner's final bid as an exact cost observation.

The other diagnostics point the same way. The auction-level heterogeneity
correction removes common scale shocks before simulation. A Convite first-price
GPV recovery and the Preg\~ao drop-out recovery line up in the key
pharmaceutical non-SME pre-period cell after that correction. A Gaussian-copula
relaxation allows within-auction cost correlation up to $\rho_c=\ipvRhoMax$;
across that grid, the within-auction share moves by less than
\ipvShareDriftPp{} percentage points and the total effect by less than
\ipvTotalDriftPct{} percent. Tightening or loosening the
$c_\epsilon\le\filterCepsUpper$ filter and using 18-, 12-, or 6-month windows
also leave the within-auction component as the larger component in both
classes. Online Appendix~OA-C reports the cost-distribution quantiles and
primitive-stability diagnostics behind these summaries.

\subsection{Protected-pool composition}

The main modeling threat is that $S_3-S_2$ mixes additional SME bidders with a
changed post-policy SME cost distribution. The strict-invariance benchmark
sets $F_c^{\mathrm{SME,Post}}=F_c^{\mathrm{SME,Pre}}$ while keeping observed
post-policy SME entry counts. The total price effect remains positive:
\siDeltaTotalNp{} in non-pharmaceuticals and \siDeltaTotalPh{} in
pharmaceuticals. The within-auction shares rise to \strictShareNp{} and
\strictSharePh{} percent. Composition therefore matters, but it is not what
makes the exclusion component dominate the price decomposition.

Composition does matter for the policy ranking in pharmaceuticals.
Post-policy pharmaceutical SME participation is heavily composed of new firms:
\turnoverPhPctNewFirms{} percent of post-period pharma SME firms are absent
from the pre-period pool, and they account for \turnoverPhPctNewBids{} percent
of post-period pharma SME bids. In non-pharmaceuticals, new firms account for
\turnoverNpPctNewBids{} percent of post-period SME bids. The pharmaceutical
welfare ranking is therefore conditional, while the non-pharmaceutical ranking
is the more stable policy result.

\subsection{Coordination}

Coordination is the conduct threat that would most directly contaminate the
drop-out interpretation. The data show bidder clustering: Conley-style
close-pair screens \citep{conley2016} reject the independent-draw null in all
four class-by-period cells, and persistent-pair diagnostics in the spirit of
\citet{bajariye2003} flag excess repeated co-bidding. The relevant question is
whether clustering increases after non-SMEs are removed. It does not. The
Conley realized close-pair share is essentially flat in non-pharmaceuticals
(\collConleyRealPreNp{} to \collConleyRealPostNp{} percent) and falls in
pharmaceuticals (\collConleyRealPrePh{} to \collConleyRealPostPh{} percent).
The Bajari--Ye $T_1$ observed-to-null ratio also falls in both classes. Residual
baseline clustering remains a limitation, but the strongest alternative story,
a post-restriction coordination shock among surviving SMEs, is not supported by
the diagnostics. Online Appendix~OA-D reports the coordination-screen table.

\subsection{Policy benchmark}

The price-preference comparison is a static design benchmark, so the relevant
robustness checks are mainly mechanical. The non-pharmaceutical ranking remains
$V_3 \succ V_0$ over $\lambda\in[\welfLambdaRangeLo,\welfLambdaRangeHi]$ under
both the main and strict-invariance specifications; the pharmaceutical ranking
remains conditional on the treatment of the post-policy SME pool. The result is
also not an artifact of exact Vickrey equivalence. A Maskin--Riley first-price
upper-bound adjustment \citep{maskinriley2000} is at most
\maskinRileyBoundPct{} percent of the reference price, while the simulated
$V_0$--$V_3$ price gap is \vZeroVThreeGapLo--\vZeroVThreeGapHi{} percent.
Online Appendix~OA-E reports the preference grid and welfare-weight
sensitivities.

The checks leave a narrow claim. The result does not depend on the reduced-form
DiD carrying the whole design, on a single treatment of winner censoring, or on
holding the post-policy SME pool fixed. When the excluded bidders are the ones
disciplining the order statistic, set-asides are costly because they replace
competition with eligibility. The protected-pool response offsets part of that
cost, but it does not replace the competitive discipline that the policy
removes.
